\documentclass[12pt]{article}

\usepackage[a4paper,margin=1in]{geometry}
\usepackage[T1]{fontenc}
\usepackage[utf8]{inputenc}
\usepackage{lmodern}
\usepackage{amsmath,amssymb}
\usepackage{booktabs}
\usepackage{longtable}
\usepackage{array}
\usepackage{tabularx}
\usepackage{enumitem}
\usepackage{hyperref}
\usepackage{setspace}
\usepackage{ragged2e}

\setstretch{1.08}
\setlength{\parskip}{0.5em}
\setlength{\parindent}{0pt}

\hypersetup{
	colorlinks=true,
	linkcolor=black,
	citecolor=black,
	urlcolor=blue,
	pdftitle={Collected Results and Proof Status Note for the General Theory of Cognitive Structuring},
	pdfauthor={Kostiantyn Osmolovskyi}
}

\setlist[itemize]{topsep=3pt,itemsep=2pt,parsep=0pt}
\setlist[enumerate]{topsep=3pt,itemsep=2pt,parsep=0pt}


\title{Collected Results and Proof Status Note\\
	\small for the General Theory of Cognitive Structuring}

\author{
	Kostiantyn Osmolovskyi\\
	\small Independent Researcher, Ukraine\\
	\small ORCID: 0009-0006-3144-7237\\
	\small \texttt{constantinosmol@gmail.com}
}

\date{}

% ---------------------------------------------------------------------------
% Part of a series: General Theory of Cognitive Structuring (GTCS)
% Autor: Kostiantyn Osmolovskyi (ORCID: 0009-0006-3144-7237)
% License: Creative Commons Attribution 4.0 International (CC BY 4.0)
% Repository: https://github.com/k-osmolovskyi/k-osmolovskyi.github.io
% DOI: https://doi.org/10.5281/zenodo.19705594
% ---------------------------------------------------------------------------


\begin{document}
	\maketitle

\begin{center}
	\small
	Verification Report No. 2026-1\\
	Version 1.0
\end{center}

\vspace{0.5em}	
	
	\section{Purpose of this note}
	
	This document provides a collected and normalized register of the main results of the
	\emph{General Theory of Cognitive Structuring}. Its purpose is not to replace the primary
	papers, nor to function as a full proof compendium, but to establish a canonical author-side layer answering the following question:
	
	\begin{quote}
		What exactly has been established in the series, in what form, and with what proof status?
	\end{quote}
	
	The note is intended to support external verification, later proof consolidation, and
	future development of a full proof compendium. It collects the main results of the series in a normalized format and distinguishes:
	\begin{itemize}[leftmargin=2em]
		\item definitional items,
		\item operators and constructions,
		\item architectural claims,
		\item proposition-level and theorem-level results,
		\item proof-complete results,
		\item proof-sketch results,
		\item and results that remain open for later consolidation.
	\end{itemize}
	
	\section{Scope and limits}
	
	This note is deliberately intermediate in scope.
	
	It does \emph{not} attempt to:
	\begin{itemize}[leftmargin=2em]
		\item reproduce all proofs in full;
		\item replace the original papers;
		\item function as a complete theorem numbering system for the entire series;
		\item decide every local notational or editorial variant;
		\item present every minor consequence or corollary.
	\end{itemize}
	
	Its function is narrower:
	\begin{itemize}[leftmargin=2em]
		\item to collect the main results of the series in one place;
		\item to normalize their wording for cross-paper verification;
		\item to assign explicit proof status to each result;
		\item to identify which results are already proof-complete and which still require later consolidation.
	\end{itemize}
	
	\section{Proof status categories}
	
	The following proof status categories are used throughout this note.
	
	\begin{longtable}{>{\RaggedRight\arraybackslash}p{0.10\textwidth}
			>{\RaggedRight\arraybackslash}p{0.75\textwidth}}
		\toprule
		\textbf{Code} & \textbf{Meaning} \\
		\midrule
		\endfirsthead
		\toprule
		\textbf{Code} & \textbf{Meaning} \\
		\midrule
		\endhead
		
		D & Definitional item. No proof is required because the item functions as a definition, operator, or construction rather than as a derived result. \\
		
		A & Architectural claim. The item is part of the intended architectural interpretation of the framework but is not presented as a theorem-grade result in canonical form. \\
		
		S & Proof sketch in source. The result is stated and supported in the source paper, but the proof remains sketch-level or relies on compressed exposition. \\
		
		P & Full proof in source. The result is stated and proved in the source paper in a form sufficient for direct canonical reference. \\
		
		C & Canonically normalized here. The result is collected and normalized in this note from one or more source papers, but this note does not itself upgrade a sketch into a full proof. \\
		
		O & Open or pending consolidation. The result is part of the intended framework but still requires later proof normalization, clarification, or formal extraction for canonical proof-level use. \\
		
		\bottomrule
	\end{longtable}
	
	\section{Reading rules for collected results}
	
	The following rules govern the present note.
	
	\begin{enumerate}[leftmargin=2em]
		\item This note does not introduce new theoretical results.
		\item Normalization of wording does not broaden the scope of the original source.
		\item A result marked \textbf{S} remains sketch-level even if its wording here is cleaner than in the source paper.
		\item A result marked \textbf{P} is treated as proof-complete in its source, but not necessarily re-proved in full here.
		\item Later synthetic papers do not retroactively upgrade the proof status of earlier papers unless this is explicitly stated.
		\item If a result depends on assumptions local to a particular paper, that locality should be preserved in its collected form.
		\item If a result appears in more than one paper, priority is given to:
		\begin{itemize}[leftmargin=2em]
			\item earliest explicit statement,
			\item strongest primary formalization,
			\item and narrowest non-overclaiming formulation.
		\end{itemize}
	\end{enumerate}
	
	\section{Result classification scheme}
	
	For organizational purposes, the collected results are grouped into the following families:
	
	\begin{itemize}[leftmargin=2em]
		\item \textbf{R0} Definitions and operators
		\item \textbf{R1} Geometric results
		\item \textbf{R2} Regulatory and overload results
		\item \textbf{R3} Invariant and compression results
		\item \textbf{R4} Representation and accessibility results
		\item \textbf{R5} Identity and multi-system results
		\item \textbf{R6} Synthetic framework-level results
	\end{itemize}
	
	This grouping is organizational and does not replace paper order. Its purpose is to make later
	proof consolidation easier.
	
	\section{Master table of collected results}
	
	\begin{longtable}{>{\RaggedRight\arraybackslash}p{0.06\textwidth}
			>{\RaggedRight\arraybackslash}p{0.22\textwidth}
			>{\RaggedRight\arraybackslash}p{0.12\textwidth}
			>{\RaggedRight\arraybackslash}p{0.10\textwidth}
			>{\RaggedRight\arraybackslash}p{0.14\textwidth}
			>{\RaggedRight\arraybackslash}p{0.14\textwidth}
			>{\RaggedRight\arraybackslash}p{0.12\textwidth}}
		\toprule
		\textbf{Result ID} & \textbf{Result title} & \textbf{Type} & \textbf{Proof status} & \textbf{Main source} & \textbf{Depends on} & \textbf{Assumptions} \\
		\midrule
		\endfirsthead
		\toprule
		\textbf{Result ID} & \textbf{Result title} & \textbf{Type} & \textbf{Proof status} & \textbf{Main source} & \textbf{Depends on} & \textbf{Assumptions} \\
		\midrule
		\endhead
		
		R0.01 & Admissibility of structural updating & Definition & D & entry paper & processing /updating distinction & --- \\
		
		R1.01 & Coherence as distance to stability region & Definition /Formal & P & coherence paper & configuration space; stability region & metric structure \\
		
		R2.01 & Admissibility depends on joint regulatory state & Formal & P & structural admissibility & tension; overload memory; admissibility operator & admissibility framework \\
		
		R2.02 & Overload memory is not reducible to instantaneous tension & Proposition & P & overload formation & overload memory; path dependence & overload dynamics \\
		
		R2.03 & Regulation is trajectory-dependent & Proposition & P & trajectory regulation & overload memory; admissibility operator & boundary deformation / path dependence \\
		
		R3.01 & Invariants are historically constituted structural constraints & Definition & D & invariants paper & architecture /process distinction & --- \\
		
		R3.02 & Invariant architecture induces stability geometry & Theorem /Proposition & P & invariants paper & invariants; admissibility profiles & regularity / continuity \\
		
		R3.03 & Compression is not equivalent to simplification & Proposition & P & compression paper & compression criterion; expected overload & compression framework \\
		
		R3.04 & Compression generates a new invariant & Proposition & P & compression paper & compression operation; invariant concept & compression framework \\
		
		R4.01 & Coherence representation is an internal regulatory variable, not full geometric reconstruction & Definition /Proposition & P & emergence paper & partial observability; local signals & representation framework \\
		
		R4.02 & Pre-symbolic admissibility is a distinct prior filtering layer & Definition & D & pre-symbolic paper & signal domain; regulatory layering & --- \\
		
		R4.03 & Restricted accessibility of coherence & Proposition & P & restricted accessibility paper & geometric coherence; admissible discrepancy domain & accessibility framework \\
		
		R4.04 & Non-injectivity of regulatory accessibility & Proposition & P/S & restricted accessibility paper & restricted accessibility; admitted-signal structure & accessibility framework \\
		
		R5.01 & Identity as regulatory attractor & Definition /Architectural claim & S & identity paper & overload; trajectory dynamics; bounded regulation & long-run dynamic assumptions \\
		
		R5.02 & Inter-system conflict as incompatibility of admissible discrepancy domains & Definition /Proposition & P/S & inter-system conflict paper & admissible discrepancy domains; compatibility criterion & multi-system setting \\
		
		R6.01 & Layered admissibility architecture of the full framework & Synthetic & C & general theory & results from R1--R5 & synthetic integration \\
		
		\bottomrule
	\end{longtable}
	
	\section{Normalized statements}
	
	\subsection{R0. Definitions and operators}
	
	\subsubsection*{R0.01. Admissibility of structural updating}
	\textbf{Type:} Definition \\
	\textbf{Statement:} Admissibility of structural updating refers to the condition under which
	already stabilized organization is available for structural updating, rather than merely
	subject to processing, evaluation, or local adjustment. \\
	\textbf{Dependencies:} processing/updating distinction. \\
	\textbf{Assumptions:} none beyond the entry-level architectural distinction. \\
	\textbf{Source:} \emph{Structural Updating and the Limits of Cognitive Change}. \\
	\textbf{Proof status:} D. \\
	\textbf{Note:} This item functions as an entry distinction rather than a derived result.
	
	\subsection{R1. Geometric results}
	
	\subsubsection*{R1.01. Coherence as distance to stability region}
	\textbf{Type:} Definition/Formal result \\
	\textbf{Statement:} Coherence is defined geometrically as the distance between the current
	configuration $x_t$ and the relevant admissible stability region, in later grounded form
	$U_X(I_t)$. \\
	\textbf{Dependencies:} configuration space; current configuration; stability region. \\
	\textbf{Assumptions:} metric or admissible distance structure over the configuration space. \\
	\textbf{Source:} \emph{Coherence Evaluation in Cognitive Systems}. \\
	\textbf{Proof status:} P. \\
	\textbf{Note:} The later invariant-grounded notation strengthens the architectural grounding
	of the earlier stability-region formulation.
	
	\subsection{R2. Regulatory and overload results}
	
	\subsubsection*{R2.01. Admissibility depends on joint regulatory state}
	\textbf{Type:} Formal result \\
	\textbf{Statement:} Structural admissibility of updating is not determined by instantaneous
	tension alone, but by the joint regulatory state including tension and overload memory. \\
	\textbf{Dependencies:} structural tension; overload memory; admissibility operator. \\
	\textbf{Assumptions:} admissibility operator defined on joint regulatory state. \\
	\textbf{Source:} \emph{Structural Admissibility in Cognitive Systems}. \\
	\textbf{Proof status:} P. \\
	\textbf{Note:} This is one of the central formal claims of the regulatory branch.
	
	\subsubsection*{R2.02. Overload memory is not reducible to instantaneous tension}
	\textbf{Type:} Proposition \\
	\textbf{Statement:} There is no function of instantaneous tension alone that captures overload
	memory for all admissible trajectories of the system. \\
	\textbf{Dependencies:} overload memory definition; excess accumulation; path dependence. \\
	\textbf{Assumptions:} nontrivial trajectory variation; retained overload state. \\
	\textbf{Source:} \emph{Overload Formation in Cognitive Processing}. \\
	\textbf{Proof status:} P. \\
	\textbf{Note:} This result supports the claim that overload contributes a genuine additional
	state dimension.
	
	\subsubsection*{R2.03. Regulation is trajectory-dependent}
	\textbf{Type:} Proposition \\
	\textbf{Statement:} Two systems with identical instantaneous tension may differ in regulatory
	state if their overload histories differ. \\
	\textbf{Dependencies:} overload memory; admissibility operator; stability region in regulatory
	phase space. \\
	\textbf{Assumptions:} admissibility depends jointly on present tension and overload memory. \\
	\textbf{Source:} \emph{Trajectory-Dependent Regulation in Cognitive Systems}. \\
	\textbf{Proof status:} P. \\
	\textbf{Note:} This result is the main trajectory-level consequence of overload retention.
	
	\subsection{R3. Invariant and compression results}
	
	\subsubsection*{R3.01. Invariants are historically constituted structural constraints}
	\textbf{Type:} Definition \\
	\textbf{Statement:} Invariants are structural constraints that participate in processing while
	being preserved under ordinary processing dynamics and arising from prior stabilization or
	compression. \\
	\textbf{Dependencies:} architecture/process distinction. \\
	\textbf{Assumptions:} none beyond the architectural ontology of the series. \\
	\textbf{Source:} \emph{Invariants in Cognitive Architectures}. \\
	\textbf{Proof status:} D. \\
	\textbf{Note:} This is a foundational definitional node for the structural architecture branch.
	
	\subsubsection*{R3.02. Invariant architecture induces stability geometry}
	\textbf{Type:} Theorem/Proposition \\
	\textbf{Statement:} Under appropriate regularity conditions, invariant structure induces a
	well-posed region of structurally stable configurations and thereby grounds geometric coherence. \\
	\textbf{Dependencies:} invariants; admissibility profiles; thresholds. \\
	\textbf{Assumptions:} continuity / regularity assumptions on admissibility profiles and
	configuration space. \\
	\textbf{Source:} \emph{Invariants in Cognitive Architectures}. \\
	\textbf{Proof status:} P. \\
	\textbf{Note:} This result strengthens the architectural grounding of the earlier geometric branch.
	
	\subsubsection*{R3.03. Compression is not equivalent to simplification}
	\textbf{Type:} Proposition \\
	\textbf{Statement:} Structural compression and structural simplification are not equivalent,
	because reduction in structural size does not by itself guarantee reduction of expected overload. \\
	\textbf{Dependencies:} admissible updating; expected overload criterion; invariant structure. \\
	\textbf{Assumptions:} compression framework. \\
	\textbf{Source:} \emph{Structural Compression in Cognitive Architectures}. \\
	\textbf{Proof status:} P. \\
	\textbf{Note:} This result blocks a major category error in later architectural reading.
	
	\subsubsection*{R3.04. Compression generates a new invariant}
	\textbf{Type:} Proposition \\
	\textbf{Statement:} If an admissible structural update is a compression, then the compressed
	representation enters the updated architecture as a new invariant. \\
	\textbf{Dependencies:} compression operation; invariant definition. \\
	\textbf{Assumptions:} compression framework. \\
	\textbf{Source:} \emph{Structural Compression in Cognitive Architectures}. \\
	\textbf{Proof status:} P. \\
	\textbf{Note:} This result links architectural evolution and invariant formation.
	
	\subsection{R4. Representation and accessibility results}
	
	\subsubsection*{R4.01. Coherence representation is an internal regulatory variable, not full geometric reconstruction}
	\textbf{Type:} Definition/Proposition \\
	\textbf{Statement:} Coherence representation is an internal regulatory variable constructed from
	available structural signals and does not reconstruct the full geometry of coherence. \\
	\textbf{Dependencies:} geometric coherence; partial observability; local signals. \\
	\textbf{Assumptions:} representation framework under bounded observability. \\
	\textbf{Source:} \emph{Emergence of Coherence Representation}. \\
	\textbf{Proof status:} P. \\
	\textbf{Note:} This is the central entry result of the representational branch.
	
	\subsubsection*{R4.02. Pre-symbolic admissibility is a distinct prior filtering layer}
	\textbf{Type:} Definition \\
	\textbf{Statement:} Pre-symbolic admissibility is a distinct regulatory layer operating before
	stable representation and before structural updating. \\
	\textbf{Dependencies:} signal domain; regulatory layering. \\
	\textbf{Assumptions:} none beyond the access-layer architecture. \\
	\textbf{Source:} \emph{Pre-Symbolic Admissibility in Cognitive Systems}. \\
	\textbf{Proof status:} D. \\
	\textbf{Note:} This result anchors the access branch and prevents conflation of operators.
	
	\subsubsection*{R4.03. Restricted accessibility of coherence}
	\textbf{Type:} Proposition \\
	\textbf{Statement:} Geometric coherence may exist without becoming regulatorily accessible,
	because only admitted discrepancies can contribute to regulatory signal formation. \\
	\textbf{Dependencies:} geometric coherence; admissible discrepancy domain; regulatory accessibility. \\
	\textbf{Assumptions:} pre-symbolic filtering prior to representation. \\
	\textbf{Source:} \emph{Restricted Accessibility of Coherence in Cognitive Systems}. \\
	\textbf{Proof status:} P. \\
	\textbf{Note:} This is the main bridge result between geometry and access.
	
	\subsubsection*{R4.04. Non-injectivity of regulatory accessibility}
	\textbf{Type:} Proposition \\
	\textbf{Statement:} Distinct global configurations may become indistinguishable at the level of
	regulatorily accessible signals when the differences between them do not enter the admissible
	discrepancy domain. \\
	\textbf{Dependencies:} restricted accessibility; admissible discrepancy structure. \\
	\textbf{Assumptions:} access-layer filtering and signal construction. \\
	\textbf{Source:} \emph{Restricted Accessibility of Coherence in Cognitive Systems}. \\
	\textbf{Proof status:} S/P. \\
	\textbf{Note:} This item may later benefit from more explicit canonical extraction.
	
	\subsection{R5. Identity and multi-system results}
	
	\subsubsection*{R5.01. Identity as regulatory attractor}
	\textbf{Type:} Definition/Architectural claim \\
	\textbf{Statement:} Identity is not introduced as a primitive representational object, but is
	derived as a structural consequence of bounded regulation under overload-sensitive dynamics. \\
	\textbf{Dependencies:} overload; trajectory dependence; admissibility boundary. \\
	\textbf{Assumptions:} long-run regulatory setting. \\
	\textbf{Source:} \emph{Identity as a Regulatory Attractor}. \\
	\textbf{Proof status:} S. \\
	\textbf{Note:} The result is central, but some long-run assumptions remain local to this branch.
	
	\subsubsection*{R5.02. Inter-system conflict as incompatibility of admissible discrepancy domains}
	\textbf{Type:} Definition/Proposition \\
	\textbf{Statement:} Inter-system conflict is defined not as disagreement of representations, but
	as the absence of a shared admissible discrepancy structure sufficient for coordinated regulation. \\
	\textbf{Dependencies:} admissible discrepancy domains; compatibility criterion; multi-system setting. \\
	\textbf{Assumptions:} admissibility-constrained inter-system interaction. \\
	\textbf{Source:} \emph{Inter-System Conflict Geometry in Cognitive Systems}. \\
	\textbf{Proof status:} S/P. \\
	\textbf{Note:} This is the defining result of the inter-system conflict branch.
	
	\subsection{R6. Synthetic framework-level results}
	
	\subsubsection*{R6.01. Layered admissibility architecture of the full framework}
	\textbf{Type:} Synthetic result \\
	\textbf{Statement:} The series as a whole is organized by layered admissibility, including the
	distinction between structural inconsistency, admissibility of discrepancy, representational
	organization, and admissibility of structural updating. \\
	\textbf{Dependencies:} results from geometric, regulatory, invariant, representational, and
	access branches. \\
	\textbf{Assumptions:} synthetic integration across the series. \\
	\textbf{Source:} \emph{General Theory of Cognitive Structuring}. \\
	\textbf{Proof status:} C. \\
	\textbf{Note:} This is a synthetic convergence result rather than a new theorem-grade source of
	the earlier branches.
	
	\section{Proof status by category}
	
	\subsection{Results with full proof in source}
	
	The following core results are currently treated as proof-complete in their source papers:
	\begin{itemize}[leftmargin=2em]
		\item R1.01
		\item R2.01
		\item R2.02
		\item R2.03
		\item R3.02
		\item R3.03
		\item R3.04
		\item R4.01
		\item R4.03
	\end{itemize}
	
	\subsection{Results stated with proof sketch or compressed support}
	
	The following items remain at least partly sketch-level in their current collected form:
	\begin{itemize}[leftmargin=2em]
		\item R4.04
		\item R5.01
		\item R5.02
	\end{itemize}
	
	\subsection{Architectural claims not yet collected as theorem-grade results}
	
	The following items are central to the framework but remain architectural or synthetic in status:
	\begin{itemize}[leftmargin=2em]
		\item R0.01
		\item R3.01
		\item R4.02
		\item R6.01
	\end{itemize}
	
	\subsection{Items likely to require later consolidation}
	
	The following families are likely candidates for later proof-compendium extraction:
	\begin{itemize}[leftmargin=2em]
		\item theorem-grade long-run results in the identity branch;
		\item full canonical normalization of non-injectivity results in the accessibility branch;
		\item collected inter-system result family with explicit local assumptions;
		\item synthetic framework-level claims that should remain explicitly non-theorem-grade.
	\end{itemize}
	
	\section{Cross-paper normalization note}
	
	This note follows the normalization principles below.
	
	\begin{enumerate}[leftmargin=2em]
		\item Where a notion appears first in operational form and later in stronger formalized form,
		the collected statement preserves the stronger form while recording source priority.
		\item Earlier notation such as $U$ may later appear as $U_X(I_t)$; the later form is treated
		as architectural strengthening rather than as a distinct unrelated object.
		\item Representational variables such as $\widehat{C}_t$ or related forms are normalized as
		internal regulatory variables and never substituted for geometric coherence itself.
		\item Synthetic papers may collect earlier results but do not retroactively become the source
		of those results.
	\end{enumerate}
	
	\section{Priority results for future proof compendium}
	
	The following groups are the most natural candidates for future canonical extraction into a
	full proof compendium.
	
	\begin{enumerate}[leftmargin=2em]
		\item \textbf{Core geometric and regulatory results}
		\begin{itemize}[leftmargin=2em]
			\item R1.01
			\item R2.01
			\item R2.02
			\item R2.03
		\end{itemize}
		
		\item \textbf{Invariant and compression results}
		\begin{itemize}[leftmargin=2em]
			\item R3.02
			\item R3.03
			\item R3.04
		\end{itemize}
		
		\item \textbf{Representation and accessibility results}
		\begin{itemize}[leftmargin=2em]
			\item R4.01
			\item R4.03
			\item R4.04
		\end{itemize}
		
		\item \textbf{Identity and multi-system branch}
		\begin{itemize}[leftmargin=2em]
			\item R5.01
			\item R5.02
		\end{itemize}
	\end{enumerate}
	
	These are the most likely results to benefit from full normalized proof presentation in a later
	canonical proof layer.
	
	\section{Tables for compact use}
	
	\subsection*{Table A. Proof status legend}
	
	\begin{longtable}{>{\RaggedRight\arraybackslash}p{0.10\textwidth}
			>{\RaggedRight\arraybackslash}p{0.75\textwidth}}
		\toprule
		\textbf{Code} & \textbf{Meaning} \\
		\midrule
		D & Definitional item; no proof required. \\
		A & Architectural claim. \\
		S & Proof sketch in source. \\
		P & Full proof in source. \\
		C & Canonically normalized here. \\
		O & Open / pending consolidation. \\
		\bottomrule
	\end{longtable}
	
	\subsection*{Table B. Compact result index}
	
	\begin{longtable}{>{\RaggedRight\arraybackslash}p{0.12\textwidth}
			>{\RaggedRight\arraybackslash}p{0.48\textwidth}
			>{\RaggedRight\arraybackslash}p{0.12\textwidth}
			>{\RaggedRight\arraybackslash}p{0.18\textwidth}}
		\toprule
		\textbf{ID} & \textbf{Title} & \textbf{Status} & \textbf{Source} \\
		\midrule
		\endfirsthead
		\toprule
		\textbf{ID} & \textbf{Title} & \textbf{Status} & \textbf{Source} \\
		\midrule
		\endhead
		
		R1.01 & Coherence as distance to stability region & P & coherence paper \\
		R2.01 & Admissibility depends on joint regulatory state & P & structural admissibility \\
		R2.02 & Overload memory is not reducible to instantaneous tension & P & overload formation \\
		R2.03 & Regulation is trajectory-dependent & P & trajectory regulation \\
		R3.02 & Invariant architecture induces stability geometry & P & invariants paper \\
		R3.03 & Compression is not equivalent to simplification & P & compression paper \\
		R3.04 & Compression generates a new invariant & P & compression paper \\
		R4.01 & Coherence representation is not full geometric reconstruction & P & emergence paper \\
		R4.03 & Restricted accessibility of coherence & P & restricted accessibility \\
		R5.01 & Identity as regulatory attractor & S & identity paper \\
		R5.02 & Inter-system conflict as incompatibility of admissible discrepancy domains & S/P & conflict paper \\
		R6.01 & Layered admissibility architecture of the framework & C & general theory \\
		\bottomrule
	\end{longtable}
	
	\section{Closing note}
	
	A paper series may remain difficult to verify even when its formal content is substantial, if its
	main results remain distributed across multiple documents and levels of analysis. The purpose of
	the present note is to reduce that difficulty by providing a canonical author-side collected layer
	of results together with proof status.
	
	This document is not the final proof layer of the \emph{General Theory of Cognitive Structuring}.
	It is a controlled intermediate step toward such a layer.
	
\end{document}